\documentclass[11pt,a4paper]{article}
\usepackage[T1]{fontenc}
\usepackage[utf8]{inputenc}
\usepackage{lmodern}
\usepackage[margin=2.5cm]{geometry}
\usepackage{booktabs}
\usepackage{tabularx}
\usepackage{enumitem}
\usepackage{listings}
\usepackage{xcolor}
\usepackage[hidelinks]{hyperref}

\lstset{
  basicstyle=\ttfamily\small,
  columns=fullflexible,
  keepspaces=true,
  frame=single,
  rulecolor=\color{gray!50},
  backgroundcolor=\color{gray!8},
  breaklines=true,
  upquote=true,
  showstringspaces=false,
  xleftmargin=0.5em,
  xrightmargin=0.5em
}
\setlength{\emergencystretch}{3em}
\newcommand{\kbd}[1]{\texttt{#1}}
\newcolumntype{L}{>{\raggedright\arraybackslash}X}

\title{LeRobot Installation Manual\\[0.3em]\large Installing LeRobot on another Mac}
\author{}
\date{8 November 2026}

\begin{document}
\maketitle

This manual explains how to install LeRobot, the OCaml drawing robot, on a
Mac where it has never been installed. Once it is done, use the
\emph{LeRobot User Manual} (\kbd{LeRobot\_Manual.pdf}) to draw and to run
robot programs.

The installation takes about half an hour, mostly spent waiting for
downloads. You need:
\begin{itemize}
\item a Mac with a recent macOS (Apple Silicon or Intel);
\item an administrator account: several steps ask for the Mac password;
\item an Internet connection;
\item a copy of the \kbd{LeRobot} folder (USB key, AirDrop, cloud drive, \dots).
\end{itemize}

All commands are typed in the \textbf{Terminal} application
(\emph{Applications} $\rightarrow$ \emph{Utilities} $\rightarrow$
\emph{Terminal}). Copy each command in one piece, paste it into the
Terminal, and press Return. When \kbd{sudo} asks for a password, type the
Mac password: nothing is shown while you type, which is normal.

\tableofcontents

%======================================================================
\section{Apple's command line tools}

These provide the C compiler and \kbd{make} that OCaml needs.
\begin{lstlisting}[language=bash]
xcode-select --install
\end{lstlisting}
Click \emph{Install} in the window that appears and wait for it to finish.
If the Terminal answers that the tools are already installed, go on to the
next step.

%======================================================================
\section{Homebrew}

Homebrew installs the other programs. Skip this step if
\kbd{brew --version} already prints a version number. Otherwise run the
installer (the same command is given on \url{https://brew.sh}, where it can
also be copied):
\begin{lstlisting}[language=bash]
/bin/bash -c \
  "$(curl -fsSL https://raw.githubusercontent.com/Homebrew/install/HEAD/install.sh)"
\end{lstlisting}
At the end, the installer prints a few \emph{Next steps} commands. Run them:
they make the \kbd{brew} command available in every new Terminal window. On
an Apple Silicon Mac they amount to:
\begin{lstlisting}[language=bash]
echo 'eval "$(/opt/homebrew/bin/brew shellenv)"' >> ~/.zprofile
eval "$(/opt/homebrew/bin/brew shellenv)"
\end{lstlisting}

%======================================================================
\section{XQuartz}

XQuartz is the display server in which the robot window opens.
\begin{lstlisting}[language=bash]
brew install --cask xquartz
\end{lstlisting}
Then \textbf{log out of the Mac and log in again} (Apple menu
$\rightarrow$ \emph{Log Out}). XQuartz only works properly after a new
session.

%======================================================================
\section{OCaml, dune and the graphics library}

\begin{enumerate}
\item Install opam, the OCaml package manager:
\begin{lstlisting}[language=bash]
brew install opam
\end{lstlisting}
\item Initialise it. This downloads and compiles OCaml, which takes several
  minutes. Answer \kbd{y} if it asks whether to modify the shell
  configuration.
\begin{lstlisting}[language=bash]
opam init -y && eval $(opam env)
\end{lstlisting}
\item Install the build system and the libraries LeRobot uses:
\begin{lstlisting}[language=bash]
opam install -y dune graphics ocamlfind && eval $(opam env)
\end{lstlisting}
\end{enumerate}
Check: \kbd{ocaml -version} and \kbd{dune --version} each print a version
number.

%======================================================================
\section{The LeRobot folder}

\begin{enumerate}
\item Copy the \kbd{LeRobot} folder anywhere, for example onto the
  \textbf{Desktop}. The OCaml programs are in its folder
  \kbd{LeRobot\_OCaml\_FR}; the \kbd{LeRobotOCaml.app} icon must stay in that
  folder, next to \kbd{robot.sh}. The whole \kbd{LeRobot} folder can be
  moved or renamed later: nothing in it depends on where it is.

  The \kbd{LeRobot} folder can be anywhere (on the
Desktop, in Documents, \dots): everything in it is found relative to it.
The commands of this manual are typed in a Terminal window \emph{opened in
the \kbd{LeRobot} folder}: type \kbd{cd} and a space, drag the
\kbd{LeRobot} folder from the Finder onto the Terminal window, and press
Return (or, in the Finder, right-click the folder and choose
\emph{Services} $\rightarrow$ \emph{New Terminal at Folder}).
\item Compile it:
\begin{lstlisting}[language=bash]
cd LeRobot_OCaml_FR && eval $(opam env) && dune build
\end{lstlisting}
  If the folder was copied with its \kbd{\_build} subfolder from another Mac
  and the build fails, delete it with \kbd{dune clean} and build again.
\end{enumerate}

%======================================================================
\section{Start XQuartz automatically after each restart}

At every restart, macOS empties the folder \kbd{/tmp}, which removes the
folder \kbd{/tmp/.X11-unix} that XQuartz needs. XQuartz cannot recreate it
itself, because it must belong to the system administrator (root). Without
it, XQuartz opens and closes at once and the robot window cannot appear.

The file \kbd{org.xquartz.x11-unix.plist}, in the \kbd{LeRobot} folder,
is a small startup task that recreates this folder at every restart.
Install it once with these four commands. Copy the four lines together and
paste them into the Terminal; the first one asks for the Mac password.
\begin{lstlisting}[language=bash]
sudo cp org.xquartz.x11-unix.plist /Library/LaunchDaemons/
sudo chown root:wheel /Library/LaunchDaemons/org.xquartz.x11-unix.plist
sudo chmod 644 /Library/LaunchDaemons/org.xquartz.x11-unix.plist
sudo launchctl bootstrap system /Library/LaunchDaemons/org.xquartz.x11-unix.plist
\end{lstlisting}
The last command also runs the task at once, so no restart is
needed. If the \kbd{.plist} file is missing from the folder, recreate it
from Appendix~\ref{app:plist} first.

Check that it worked:
\begin{lstlisting}[language=bash]
ls -ld /tmp/.X11-unix
\end{lstlisting}
The answer must begin with \kbd{drwxrwxrwt} and contain \kbd{root wheel}.

If you ever want to remove the startup task (do not run this to install
it):
\begin{lstlisting}[language=bash]
sudo launchctl bootout system/org.xquartz.x11-unix
sudo rm /Library/LaunchDaemons/org.xquartz.x11-unix.plist
\end{lstlisting}

%======================================================================
\section{First launch}

\begin{enumerate}
\item Start XQuartz: \kbd{open -a XQuartz}.
\item Open the interactive drawing program:
\begin{lstlisting}[language=bash]
cd LeRobot_OCaml_FR
eval $(opam env)
DISPLAY=:0 dune exec ./Dessiner.exe
\end{lstlisting}
  The robot window ($800 \times 700$ pixels) opens in XQuartz. If you do not
  see it, click the XQuartz icon in the Dock. Pressing Return or clicking in
  the window ends the program.
\item \kbd{DISPLAY=:0} must always be put in front of the command: the
  default display setting may point to a stale XQuartz connection. The
  shortcut \kbd{make Dessiner}, typed in the same folder, does all of this.
\item \kbd{LeRobotOCaml.app} starts XQuartz if it is not running.
  Double-click its icon to run \kbd{Dessiner}; drop a \kbd{.ml} program on
  its icon to compile and run it.
  The first time you open \kbd{LeRobotOCaml.app}, macOS may refuse to open it because it comes from
  another Mac. Right-click the icon, choose \emph{Open}, then confirm. Allow
  it to control Terminal and to access the folder of \kbd{LeRobot}
  (Desktop, Documents, \dots) when macOS asks.
\end{enumerate}

%======================================================================
\section{Troubleshooting}

\begin{tabularx}{\textwidth}{@{}LLL@{}}
\toprule
\textbf{Symptom} & \textbf{Cause} & \textbf{Fix} \\
\midrule
\emph{command not found: brew} & Homebrew's \emph{Next steps} were not run & Run the \kbd{echo 'eval \dots'} command of Section~2, or open a new Terminal window \\
\addlinespace
\emph{command not found: dune} or \kbd{ocaml} & opam's settings not loaded in this window & Run \kbd{eval \$(opam env)} \\
\addlinespace
\kbd{opam install graphics} fails, mentioning X11 & XQuartz not installed, or the Mac not logged out and in again & Do Section~3 completely, then retry \\
\addlinespace
XQuartz opens and closes at once & \kbd{/tmp/\allowbreak .X11-unix} missing or not owned by root & Run the startup task again: \kbd{sudo launchctl kickstart system/\allowbreak org.xquartz.x11-unix}; if that fails, install it (Section~6) \\
\addlinespace
\emph{Cannot open display} & \kbd{DISPLAY=:0} missing, or XQuartz not running & Put \kbd{DISPLAY=:0} in front of the command; \kbd{open -a XQuartz} \\
\addlinespace
\emph{Boot-out failed: 3: No such process} & The removal command was run while the task was not installed & Harmless; run the installation commands of Section~6 \\
\addlinespace
\emph{Bootstrap failed: 5} & The task is already installed & Nothing to do \\
\addlinespace
XQuartz log & Detailed error messages & \kbd{\textasciitilde/\allowbreak Library/\allowbreak Logs/\allowbreak X11/\allowbreak org.xquartz.log} \\
\bottomrule
\end{tabularx}

%======================================================================
\appendix
\section{The startup task file}\label{app:plist}

If \kbd{org.xquartz.x11-unix.plist} is missing, this command recreates it
in the \kbd{LeRobot} folder (the Terminal being opened in that folder);
then install it as in
Section~6.
\begin{lstlisting}[language=bash]
cat > org.xquartz.x11-unix.plist <<'EOF'
<?xml version="1.0" encoding="UTF-8"?>
<plist version="1.0">
<dict>
  <key>Label</key>
  <string>org.xquartz.x11-unix</string>
  <key>ProgramArguments</key>
  <array>
    <string>/bin/sh</string>
    <string>-c</string>
    <string>mkdir -p /tmp/.X11-unix
chown root:wheel /tmp/.X11-unix
chmod 1777 /tmp/.X11-unix</string>
  </array>
  <key>RunAtLoad</key>
  <true/>
</dict>
</plist>
EOF
\end{lstlisting}

\end{document}
